\subsection*{A.11 source-root and angular-derivative certificate}
\label{subsec:a11-rootjet-new-20261021}
For the prescribed A.11 smooth-step derivative bumps, the true
normalized correction equations are
\[
 u(c_1+xc_2)=\Delta_I(\eta),\qquad
 2v(c_1+zc_2)+w(c_1^2+zc_2^2)=0,
\]
where $u,v,w>0$, $x=e^{2(1-\lambda)}$,
$z=e^{-2(1+2\lambda)}$ are independent of $\eta$.
Earlier source-specific bounds give
$\|B^{-1}\|_\infty<3/2$ for
$B=\left(\begin{smallmatrix}u&ux\\2v&2vz\end{smallmatrix}\right)$
and $\|Q(c)\|_\infty\le320\|c\|_\infty^2$ for
$Q(c)=(0,w(c_1^2+zc_2^2))$.

\begin{lemma}[Authentic moment-mismatch jet]
\label{lem:a11-rootjet-delta-new}
Assume the source A.10 interpolation $T_f=512$, the
pre-interpolation factorization $E(y,\eta)=c(y)(1+\eta^2)^{-1}$
and the succeeding $-\lambda$ slope plateau of length
$30\log(1/\lambda)$. For $0<\lambda\le1/100$ the
moment mismatch at the first A.11 bump center obeys
\[
 |\Delta_I|<\frac{63}{100}\lambda^{28},\qquad
 |\partial_\eta\Delta_I|<\frac14\lambda^{28}.
\]
\end{lemma}
\begin{proof}
The angular moment ratio $r=I/(XH)$ solves
$r'+(1+\ell)r=1$. It is independent of $\eta$
before A.10, hence $r_\eta=0$ upon interpolation entry.
On A.10, $\ell=-\lambda+\vartheta_f'(\log2-J_0)$,
$J_0=\log(1+\eta^2)$ and
$\vartheta_f=1-\sigma(y/T_f)$, so
$|\ell_\eta|\le|\vartheta_f'|$ and
$\int|\vartheta_f'|=1$.
The source stage bounds give $0\le r\le20/17$ and
$1+\ell\ge17/20$; differentiating the exact ODE shows
$|r_\eta(T_f)|\le20/17$.
The subsequent constant-slope interval has length
$30\log(1/\lambda)-3$ before the first bump.
Both the mismatch from its fixed point $(1-\lambda)^{-1}$
(initially of magnitude below three) and $r_\eta$
are multiplied by
$\exp[-(1-\lambda)(30\log(1/\lambda)-3)]$.
Since $30(1-\lambda)\ge29$, $e^3<21$
and $\lambda\le1/100$, one gets
$|\Delta_I|<63\lambda^{29}\le(63/100)\lambda^{28}$
and $|\Delta_{I,\eta}|<(420/17)\lambda^{29}
<(1/4)\lambda^{28}$.
\end{proof}

\begin{theorem}[Small corrected root and its first parameter jet]
\label{thm:a11-rootjet-new}
Under the lemma, the \emph{actual} nonlinear A.11 equations
have a unique solution in $\|c\|_\infty\le2\lambda^{28}$,
depending $C^1$ on $\eta$, with
\[
 \|c\|_\infty<2\lambda^{28},\qquad
 \|\partial_\eta c\|_\infty<\lambda^{28}.
\]
\end{theorem}
\begin{proof}
On the ball $R_c=2\lambda^{28}$ the map
$c\mapsto B^{-1}((\Delta_I,0)-Q(c))$ has norm at most
$\bigl((3/2)(63/100)+1920\lambda^{28}\bigr)\lambda^{28}
<R_c$ and Lipschitz constant at most
$1920\lambda^{28}<1/100$. It has a unique small fixed
point. The differentiated \emph{source} equations give
$[B+DQ(c)]c_\eta=(\Delta_{I,\eta},0)$.
The inverse norm is at most
$(3/2)/(1-1/100)$, and therefore
$\|c_\eta\|_\infty\le(25/66)\lambda^{28}
<\lambda^{28}$.
\end{proof}

Because the two source bumps are disjoint and
$\|\beta_i\|_\infty\le640/3$, this gives on their
closed convex hull (including the gap)
\[
 \left|\frac{E}{E_b}-1\right|
 \le\frac{1280}{3}\lambda^{28},\qquad
 \left|\partial_\eta(E/E_b)\right|
 \le\frac{640}{3}\lambda^{28}.
\]
It strengthens the previously assumed correction and
parameter-jet constants in the \emph{conditional}
full-hull $|N_s|/E$ theorem. It does not prove the
same profile's exact $S(\infty)=\Pi(\infty)=0$,
the future-tail integral enclosure, $Q_s$ positivity,
all exterior stress-cone stages, the Appendix B fixed
point, or the proposed unforced Theorem~1.1.

The parameter-derivative argument is essential:
$\Delta_I(\eta)=\epsilon\sin(N\eta)$ has a small
uniform value but derivative as large as $N\epsilon$.
The integral/ODE argument above, not the
small-coefficient estimate alone, supplies the
actual source's derivative bound.
